%15 min preso!
\documentclass[xcolor=table,aspectratio=169]{beamer}
\usepackage{beamerthemesplit}
\usepackage{wrapfig}
\usetheme{SPbGU}
\usepackage{pdfpages}
\usepackage{amsmath}
\usepackage{cmap}
\usepackage[T2A]{fontenc}
\usepackage[utf8]{inputenc}
\usepackage[english]{babel}
\usepackage{indentfirst}
\usepackage{amsmath}
\usepackage{tikz}
\usepackage{multirow}
\usepackage[noend]{algpseudocode}
\usepackage{algorithm}
\usepackage{algorithmicx}
\usepackage{fancyvrb}
\PassOptionsToPackage{hyphens}{url}\usepackage{hyperref} 
\usetikzlibrary{calc}
\usetikzlibrary{shapes, backgrounds}
\usetikzlibrary{arrows,automata}
\usetikzlibrary{positioning}
\usetikzlibrary{fit}
\usetikzlibrary{shapes.callouts}
\usetikzlibrary{shapes.misc}
\usepackage{xparse}
\usepackage{fontawesome}

\usepackage{etoolbox,refcount}
\usepackage{multicol}

\usepackage{tabularx}
\newcolumntype{Y}{>{\raggedleft\arraybackslash}X}

\renewcommand{\thealgorithm}{}

\newtheorem{mytheorem}{Theorem}
\renewcommand{\thealgorithm}{}

\newcommand{\tikzmark}[1]{\tikz[overlay,remember picture] \node (#1) {};}
\def\Put(#1,#2)#3{\leavevmode\makebox(0,0){\put(#1,#2){#3}}}

\newcommand{\ltz}{$< 1$}

\definecolor{links}{HTML}{2A1B81}
\hypersetup{colorlinks,linkcolor=,urlcolor=links}


\tikzset{
    state/.style={
           rectangle,
           rounded corners,
           draw=black, very thick,
           minimum height=2em,
           inner sep=2pt,
           text centered,
           },
}

\tikzset{
    invisible/.style={opacity=0,text opacity=0},
    visible on/.style={alt=#1{}{invisible}},
    alt/.code args={<#1>#2#3}{%
      \alt<#1>{\pgfkeysalso{#2}}{\pgfkeysalso{#3}} % \pgfkeysalso doesn't change the path
    },
}

\tikzset{cross/.style={cross out, draw=black, minimum size=2*(#1-\pgflinewidth), inner sep=0pt, outer sep=0pt, ultra thick},
%default radius will be 1pt. 
cross/.default={1pt}}

\NewDocumentCommand{\mycallout}{r<> O{opacity=0.8,text opacity=1} m m m}{%
\tikz[remember picture, overlay]\node[align=center, fill=cyan!20, text width=#5cm,
#2,visible on=<#1>, rounded corners,
draw,rectangle callout,anchor=pointer,callout relative pointer={(290:0.5cm)}]
at (#3) {#4};
}

\NewDocumentCommand{\mycalloutR}{r<> O{opacity=0.8,text opacity=1} m m m}{%
\tikz[remember picture, overlay]\node[align=center, fill=cyan!20, text width=#5cm,
#2,visible on=<#1>, rounded corners,
draw,rectangle callout,anchor=pointer,callout relative pointer={(30:0.8cm)}]
at (#3) {#4};
}


%callout relative pointer={(230:0.5cm)}]

\newcounter{countitems}
\newcounter{nextitemizecount}
\newcommand{\setupcountitems}{%
  \stepcounter{nextitemizecount}%
  \setcounter{countitems}{0}%
  \preto\item{\stepcounter{countitems}}%
}
\makeatletter
\newcommand{\computecountitems}{%
  \edef\@currentlabel{\number\c@countitems}%
  \label{countitems@\number\numexpr\value{nextitemizecount}-1\relax}%
}
\newcommand{\nextitemizecount}{%
  \getrefnumber{countitems@\number\c@nextitemizecount}%
}
\newcommand{\previtemizecount}{%
  \getrefnumber{countitems@\number\numexpr\value{nextitemizecount}-1\relax}%
}
\makeatother    
\newenvironment{AutoMultiColItemize}{%
\ifnumcomp{\nextitemizecount}{>}{3}{\begin{multicols}{2}}{}%
\setupcountitems\begin{itemize}}%
{\end{itemize}%
\unskip\computecountitems\ifnumcomp{\previtemizecount}{>}{3}{\end{multicols}}{}}


\beamertemplatenavigationsymbolsempty

\title[Лаборатория YADRO, Мат-Мех СПбГУ]{Лаборатория YADRO в СПбГУ}
\subtitle{на математико-механическом факультете}
\institute[СПбГУ]{
Санкт-Петербургский Государственный Университет
}

\author[Семён Григорьев]{Семён Григорьев}

\date{18 ноября 2026}


\begin{document}
{
\begin{frame}[fragile]
  \begin{table}
  \centering
  \begin{tabularx}{\linewidth}{XcX}
    \includegraphics[height=1.4cm]{pictures/YADRO_Logo.png} 
    \hfill
    & 
    & \hfill \includegraphics[height=1.6cm]{pictures/SPbGU_Logo.png}
  \end{tabularx}
  \end{table}
  \titlepage
\end{frame}
}

\begin{frame}[fragile]
  \frametitle{О лаборатории}
  \begin{itemize}
    \item Создана в 2022 году
    \item Руководитель: Яков Александрович Кириленко
    \item Объединяет преподавателей, исследователей и студентов, увлечённых системным программированием
    \vfill

    \item Перенос разностороннего опыта, навыков и знаний в индустрию
    \item Решение прикладных задач компании
    \begin{itemize}
      \item Полезных, важных, но не являющихся критичными в данный момент
    \end{itemize}
    \item Участие в работе Альянса RISC-V
    \begin{itemize}
      \item Образовательный комитет
      \item Рабочая группа <<Развитие экосистемы ПО>> Альянса RISC-V  
    \end{itemize}
  \end{itemize}
\end{frame}

\begin{frame}[fragile]
  \frametitle{Основные направления}
  \begin{minipage}{0.48\textwidth}
      \underline{\textbf{Системное программирование}}
      \small{
        \begin{itemize}\setlength\itemsep{-0.4em}
          \item Кирилл Смирнов, Дмитрий Косарев, Тимофей Фролов
          \item Ghidra, LLVM, Sail, \textbf{\href{https://github.com/Amphimixis/amphimixis}{Amphimixis}}, \textbf{\href{https://github.com/verihogg}{Verihogg}}
          \item C, C++, OCaml, ассемблер
        \end{itemize}
      }
        \underline{\textbf{Аппаратное обеспечение}}
        \small{
        \begin{itemize}\setlength\itemsep{-0.4em}
          \item Семён Григорьев, Ефим Кубышкин
          \item \href{https://github.com/Lamagraph}{LamaGraph}, \textbf{\href{https://github.com/vortexgpgpu/vortex}{Vortex GPGPU}}
          \item HDL, System Verilog, Clash, FPGA
        \end{itemize}
        }
        \underline{\textbf{\href{https://github.com/SparseLinearAlgebra}{Разреженная линейная алгебра}}}
        \small{
        \begin{itemize}\setlength\itemsep{-0.4em}
          \item Семён Григорьев, Владимир Кутуев
          \item GraphBLAS, Spla, LAGraph, CuBool
          \item C, C++, Cuda, OpenCL, GPGPU
        \end{itemize}
        }      
    \end{minipage}~  
  \begin{minipage}{0.48\textwidth}
        \underline{\textbf{Алгоритмическая статистика}}
        \small{
        \begin{itemize}\setlength\itemsep{-0.4em}
          \item Вячеслав Гориховский, Михаил Михайлов
          \item \href{https://github.com/PySATL/}{PySATL}, cеминар Statistical Software
          \item Python, математическая статистика
        \end{itemize}
        }
        \underline{\textbf{ Системы хранения данных}}
        \small{
        \begin{itemize}\setlength\itemsep{-0.4em}
          \item Анна Васенина, Алексей Дмитриевцев
          \item SPDK, ChunkFS, \textbf{\href{https://github.com/RMamonts/nfs-mamont}{NFS Mamont}}
          \item Linux Kernel, дедупликация, Rust
        \end{itemize}
        }
        \underline{\textbf{Компьютерные сети}}
        \small{
        \begin{itemize}\setlength\itemsep{-0.4em}
          \item Илья Зеленчук
          \item \href{https://github.com/mimi-net/miminet}{Miminet}, Stepik, <<Импульс>>
          \item Эмуляция сетей, VLAN, STP/RSTP, IPIP
        \end{itemize}
        }
    \end{minipage}
\end{frame}


\begin{frame}[fragile]
  \frametitle{Открытое ПО}
  \begin{itemize}
      \item Вовлечение студентов в развитие открытого ПО 
      \begin{itemize}
        \item Активное участие в международном сообществе
        \begin{itemize}
        \item Существенный вклад в резюме студента
        \end{itemize}
        \item Открытость и переиспользуемость результатов
        \item Качественная постановка экспериментов, их повторяемость (как минимум)
        \item Возможность внешней оценки и контроля
        \item Привлечение внешних участников
      \end{itemize}
      \item Лучшие инженерные практики
      \begin{itemize}
        \item У многих руководителей проектов есть актуальный опыт промышленной разработки
        \item Оформление кода и репозитория
        \item Управление процессом разработки
        \item Использование AI-агентов
        \item \ldots
      \end{itemize}
      
  \end{itemize}
\end{frame}



{
\author[Валерия Чурикова]{Валерия Чурикова}
\begin{frame}[fragile]
  \frametitle{Amphimixis: Готовность к RISC-V (Валерия Чурикова)}
  \begin{minipage}{0.7\textwidth}
  \begin{itemize}
    \item Оценка экосистемы и производительности проектов
    \item Готовность проектов к миграции
    \begin{itemize}
      \item Совместимость с \texttt{RISC-V} (RVA23)
      \item Поиск узких мест и дефектов
    \end{itemize}
    \item Методология и формальный тулинг
    \item \textbf{Amphimixis-AI: LLM-агенты}
    \begin{itemize}
      \item Интеграция агентов с \texttt{OpenCode}
      \item Пакетный режим в контейнере (\texttt{Doxis})
    \end{itemize}
    \item \textbf{17 отчётов} по открытым проектам
    \item \textbf{3 патча} в open-source: \texttt{\href{https://github.com/util-linux/util-linux/pull/4662}{util-linux}}, \texttt{\href{https://github.com/libtom/libtommath/pull/615}{libtommath}}, \texttt{\href{https://github.com/DaveGamble/cJSON/pull/1096}{cJSON}}
    \begin{itemize}
      \item Ускорение алгоритма в \texttt{util-linux} на RISC-V в \textbf{13 раз}
      \item Исправление ошибок в \texttt{libtommath} и \texttt{cJSON}
    \end{itemize}
    \item Команда: 6 человек. Руководитель: Кирилл Смирнов
  \end{itemize}
\end{minipage}
\begin{minipage}{0.26\textwidth}
  \begin{center}
    \makebox[\textwidth][c]{%
    \hspace*{-0.1cm}%
    \includegraphics[width=1.2\textwidth]{pictures/Amphimixis_Logo.png}
  }\\
    \includegraphics[width=0.6\textwidth]{pictures/Amphimixis_QR.png}
    \\[-0.2cm]
    или
    \\
    \url{amphimixis.org}
  \end{center}
\end{minipage}
\end{frame}
}

{
\author[Дмитрий Мотузов]{Дмитрий Мотузов}
\begin{frame}[fragile]
  \frametitle{Verihogg: инструменты для SystemVerilog}
  \begin{minipage}{0.7\textwidth}
  \begin{itemize}
    \item Два инструмента --- одна задача: качественный код
    \begin{itemize}
      \item \textbf{verihogg-format:} приводит код к единому стилю: отступы, переносы, выравнивание
      \item Неизвестные конструкции не трогаем: лучше пропустить, чем сломать
      \item \textbf{verihogg-lint:} статический анализ до сборки; 76 правил, 10 с автоисправлением
    \end{itemize}
    \item Открытый код; Docker-образ и сборка через Nix
    \item Руководитель: Тимофей Фролов
  \end{itemize}
\end{minipage}
\begin{minipage}{0.26\textwidth}
  \begin{center}
    \includegraphics[width=\textwidth]{pictures/sv.jpeg}\\
    \href{https://github.com/verihogg}{\includegraphics[width=0.6\textwidth]{pictures/verihogg.pdf}}\\
    или \\
    \href{https://t.me/g4rry1}{\texttt{t.me/g4rry1}}
  \end{center}
\end{minipage}
\end{frame}
}

{
\author[Алексей Дмитриевцев]{Алексей Дмитриевцев}
\begin{frame}[fragile]
  \frametitle{Marline: дедупликация по схожести для СХД}
  \begin{minipage}{0.7\textwidth}
  \begin{itemize}
    \item Content-Defined Chunking (CDC)
    \begin{itemize}
      \item Оптимизации существующих этапов пайплайна 
      \item Новые решения для интеграции с SBC
    \end{itemize}
    \item Similarity-Based Chunking (SBC)
    \begin{itemize}
      \item Похожий чанк хранится как дельта к базовому
      \item Выигрыш там, где CDC-дедупликация бессильна
    \end{itemize}
    \item Прототип дедуплицирующей ФС
    \begin{itemize}
      \item Rust, ChunkFS
      \item Скетчи, индекс похожих чанков
      \item Дельта-кодирование: Gdelta, Xdelta, Zdelta, иерархический chunk patch
    \end{itemize}
    \item[\faGears] Enterprise-сценарии
    \begin{itemize}
      \item Бэкапы: запись и чтение
    \end{itemize}
    \item Руководитель: Алексей Дмитриевцев
  \end{itemize}
\end{minipage}
\begin{minipage}{0.26\textwidth}
  \begin{center}
    \begin{tikzpicture}[font=\small, every node/.style={draw, rounded corners, minimum height=0.6cm, inner sep=3pt}]
      \node[fill=cyan!20] (a) {base};
      \node[draw=none, right=0.1cm of a] (eq) {$\approx$};
      \node[fill=orange!30, right=0.1cm of eq] (b) {chunk};
      \node[fill=green!20, below=0.45cm of eq] (d) {base + $\Delta$};
      \draw[->, thick] (b) -- (d);
    \end{tikzpicture}\\[0.3em]
    \includegraphics[width=0.6\textwidth]{pictures/marline.pdf}\\
    или \\
    \href{https://t.me/admitrievtsev}{\texttt{t.me/admitrievtsev}}
  \end{center}
\end{minipage}
\end{frame}
}

\begin{frame}[fragile]
  \frametitle{Miminet: эмулятор компьтерных сетей}
  \begin{minipage}{0.7\textwidth}
  \begin{itemize}
    \item Эмулятор компьютерных сетей, предназначенный для образовательных целей
    \begin{itemize}
      \item Помощь в самостоятельном изучении области
      \item Задачи на конструирование и настройку сетей
      \item Автоматизация проверки задач и качества усвоения материала
      \item Курсы на Stepik
      \item Отбор участников на <<Импульс>>
    \end{itemize}
    \item Эксперименты по использованию LLM в образовании
    \item Руководитель: Илья Зеленчук
  \end{itemize}
\end{minipage}
\begin{minipage}{0.26\textwidth}
  \begin{center}
    \includegraphics[width=0.7\textwidth]{pictures/miminet.png}\\
    \includegraphics[width=0.6\textwidth]{pictures/miminet.pdf}\\
    или \\
    \url{https://miminet.ru/}
  \end{center}
\end{minipage}
\end{frame}

\begin{frame}[fragile]
  \frametitle{Линейная алгебра: фундамент высокороизводительного анализа графов}
  \begin{minipage}{0.7\textwidth}
  \begin{itemize}
    \item Весь стек
    \begin{itemize}
      \item~<<Железо>>: Lamagraph, Vortex, Ventus (RISC-V GPGPU)
      \item Системное ПО: библиотеки линейной алгебры, компиляторы
      \item Прикладное ПО: анализ графов, анализ кода, базы данных
      \item Алгоритмы (разработка и анализ): графы, линейная алгебра, формальные языки, анализ сложности
    \end{itemize}
    \item Прикладные и теоретические исследования
    \begin{itemize}
      \item Вклад в международные проекты LaGraph, GraphBLAS
    \end{itemize}    
    \item  \textbf{Два постера на C++ митапе в СПб 3 декабря}
    \begin{itemize}
      \item Статический анализ кода (SVF)
      \item Анализ графов на GPU (Spla)
    \end{itemize}
    \item Руководитель: Семён Григорьев
  \end{itemize}
\end{minipage}
\begin{minipage}{0.26\textwidth}
  \begin{center}
    \includegraphics[width=0.6\textwidth]{pictures/spla.pdf}\\
    или \\
    \url{https://github.com/SparseLinearAlgebra}
  \end{center}
\end{minipage}
\end{frame}

\begin{frame}[fragile]
  \frametitle{Возможности для студентов}
  \begin{itemize}
      \item Курсовые, ВКР, магистерские, кандидатские
      \item В проектах принимают участие студенты из других вузов
      \begin{multicols}{3}
      \begin{itemize}
        \item ННГУ
        \item МФТИ
        \item ВШЭ (СПб)
        \item ИТМО
        \item ЛЭТИ
      \end{itemize}
      \end{multicols}
      \item Доклады на конференциях, профильных митапах, публикации (\textbf{с первого курса бакалавриата})
      %\begin{multicols}{2}
      \begin{itemize}
        \item {<<Инженерные системы: RISC-V \& FPGA>>, молодежная школа при <<Математическое моделирование и суперкомпьютерные технологии>>}
        \item Школа молодых учёных при форуме <<Микроэлектроника>>
        \item Современные технологии в теории и практике программирования, СПИСОК
        \item LSGDA@VLDB, GRADES-NDA@SIGMOD/PODS, SOAP@PLDI, DAMDID, PaCT, SYRCoSE
        \item Журналы различного уровня (отечественные, международные)
      \end{itemize}
      %\end{multicols}
  \end{itemize} 
\end{frame}

\begin{frame}[fragile]
  \frametitle{Заключение}
  \begin{minipage}{0.6\textwidth}
  \begin{itemize}
    \item Открытый исходный код
    \item Связь с индустрией
    \item Фундаментальные и прикладные исследования
    \item Доклады на конференциях и публикации в журналах
    \item Летняя школа и другие образовательные инициативы
  \end{itemize}
  \vspace{1cm}
\end{minipage}~
%\begin{minipage}{0.45\textwidth}
\begin{minipage}{0.35\textwidth}
 % \begin{center}
 %   \textbf{Яков Александрович Кириленко} \\
 %   \includegraphics[width=0.4\textwidth]{pictures/k.pdf}\\
 %   \textbf{!!!!!!!!!!TODO!!!!!!!!!!}
 % \end{center}
 % \end{minipage}~  
    \begin{center}
      \textbf{Семён Вячеславович Григорьев} \\
      \includegraphics[width=0.3\textwidth]{pictures/rsdpisuy.pdf}\\
      \textbf{s.v.grigoriev@mail.spbu.ru}
    \end{center}
\end{minipage}
\end{frame}


\end{document}
